\section{Plan del proyecto}

\subsection{Roles}
\begin{table}[H]
\centering
\caption{Roles del equipo CHRS}
\begin{tabularx}{\linewidth}{@{}p{4.2cm}p{3.6cm}X@{}}
\toprule
\textbf{Integrante} & \textbf{Rol} & \textbf{Responsabilidades} \\
\midrule
Silva Triviño John Jairo & Líder de Desarrollo & Repositorio, esqueleto en \texttt{/src}, modelo de comunicación, tablero e issues, video. \\
Herrera Barco Humberto Aldair & Responsable de Documentación & Documento \LaTeX, referencias, bitácora, compilación y PDF del SGA. \\
Reinoso Vélez Eduardo David & Arquitecto & Diagramas C4, justificación de la arquitectura y ADR. \\
Calderón Saltos Joseph & Responsable de Calidad & Requisitos no funcionales, CAP/PACELC, riesgos y pruebas del esqueleto. \\
\bottomrule
\end{tabularx}
\end{table}

\subsection{Tablero e issues}
El trabajo se organiza en el tablero \textit{PFC CHRS} de GitHub Projects, con las columnas
\textit{Por hacer}, \textit{En curso} y \textit{Hecho}. Cada issue tiene un responsable y la
etiqueta de su entrega (E1 o E2), y se cierra con un commit de quien lo realizó.

\subsection{Cronograma}
\begin{table}[H]
\centering
\caption{Cronograma de las entregas del PFC}
\begin{tabularx}{\linewidth}{@{}p{1.2cm}p{1.6cm}X@{}}
\toprule
\textbf{Entrega} & \textbf{Semana} & \textbf{Contenido} \\
\midrule
E1 & 5  & Definición, requisitos y arquitectura base \\
E2 & 9  & Microservicios, API y contenedores \\
E3 & 13 & Base de datos distribuida y procesamiento \\
E4 & 17 & Sistema final en capas, calidad, pruebas y defensa \\
\bottomrule
\end{tabularx}
\end{table}
